\documentclass[book.tex]{subfiles}
\begin{document}

\chapter{Parabolic Interpolation}

\label{chap:parabolic}
\noindent\rule{11cm}{0.4pt} \\
\textbf{Prerequisites:} \autoref{chap:newton_raphson} \\
\textbf{Difficulty Level:} * \\
\noindent\rule{11cm}{0.4pt}
\vspace{5mm}

In this chapter, we will learn how to use root-finding techniques to find local optima. We will understand how Newton's method for on-dimensional optimization is an extension from the Newton-Raphson method for root finding. We will also study how the parabolic interpolation method can be used for one-dimensional optimization. We then discuss how to implement the parabolic interpolation method to location the optima of a single-variable function.

\section{Parabolic Interpolation}

Parabolic interpolation method is another way of optimizing a single-variable function. It takes advantage of the fact that a second-order polynomial often provides a good approximation to the shape of $f(x)$ near an optimum. This can be visualized from the Figure \ref{fig:1}.

Just as there is only one straight line connecting two points, there is only one parabola connecting three points. Thus, if we have three points that jointly bracket an optimum, we can fit a parabola to the points. Then we can differentiate it, set the result equal to zero, and solve for an estimate of the optimal \textit{x}. It can be shown through some algebraic manipulations that the result is

\begin{align}
 x_4 = x_2 - \frac{1}{2}\dfrac{(x_2 - x_1)^2[f(x_2)-f(x_3)] - (x_2 - x_3)^2[f(x_2)-f(x_1)]}{(x_2 - x_1)[f(x_2)-f(x_3)] - (x_2 - x_3)[f(x_2)-f(x_1)]}
\end{align}

\begin{marginfigure}
	\centering
	\includegraphics[width = 2.7in]{figures/part1b/parabolic_interpolation/Parabolic.PNG}
	\centering
	\caption{Graphical Depiction of Parabolic Interpolation}
	\label{fig:1}
\end{marginfigure}

where $x_1$, $x_2$, and $x_3$ are the initial guesses, and $x_4$ is the value of x that corresponds to the optimum value of the parabolic fit to the guesses. Let's use this method to find the extremum of the same function as in Module 3.1.\\
\begin{example}
Find the minima of the unimodal function $f(x) = e^x-2x$ using parabolic interpolation method with initial guesses of $x_1 = 0$, $x_2 = 1$ and $x_3 = 2$\\
\textbf{Solution: }
\begin{fullwidth}
	\begin{Verbatim}
clear all;
f = @(x) exp(x)-2*x;
x1 = 0; x2 = 1; x3 = 2; eps = 1e-6; x_old = 100;
fx1 = f(x1); fx2 = f(x2); fx3 = f(x3);
while(1)
	num = (x2-x1)^2*(fx2-fx3)-(x2-x3)^2*(fx2-fx1);
	den = 2*((x2-x1)*(fx2-fx3)-(x2-x3)*(fx2-fx1));
	x4 = x2 - num/den;
	fx4 = f(x4);
	if fx4 < fx2
		if x4 > x2;
			x1 = x2;
			fx1 = fx2;
		else
			x3 = x2;
			fx3 = fx2;
		end
		x2 = x4;
		fx2 = fx4;
	else
		if x4 > x2;
			x3 = x4;
			fx3 = fx4;
		else
			x1 = x3;
			fx1 = fx3;
		end
	end
	if (abs((x4-x_old)/x_old) < eps)
		break
	end
	x_old = x4;
end
disp('Extremum Value')
disp(x4)
\end{Verbatim}
\end{fullwidth}
\end{example}
    
The parabolic interpolation method converges to the optimum within nine iterations to the true value of 0.6137 at x = 0.6931. The values of the guesses and their function evaluations are given below.
\begin{fullwidth}
	\begin{Verbatim}
i       x1       fx1       x2       fx2        x3        fx3       x4        fx4
1  0.000000  1.000000  1.000000  0.718282  2.000000  3.389056  0.595417  0.622953
2  0.000000  1.000000  0.595417  0.622953  1.000000  0.718282  0.662117  0.614659
3  0.595417  0.622953  0.662117  0.614659  1.000000  0.718282  0.687128  0.613742
4  0.662117  0.614659  0.687128  0.613742  1.000000  0.718282  0.691325  0.613709
5  0.687128  0.613742  0.691325  0.613709  1.000000  0.718282  0.692758  0.613706
6  0.691325  0.613709  0.692758  0.613706  1.000000  0.718282  0.693037  0.613706
7  0.692758  0.613706  0.693037  0.613706  1.000000  0.718282  0.693122  0.613706
8  0.693037  0.613706  0.693122  0.613706  1.000000  0.718282  0.693140  0.613706
9  0.693122  0.613706  0.693140  0.613706  1.000000  0.718282  0.693146  0.613706

\end{Verbatim}
\end{fullwidth}

\section{Newton's Method}

The same techniques learnt in root-finding can be extended to optimization by changing the function under consideration to $f'(x)$ from $f(x)$ i.e. the values of x where $f'(x) = 0$ are local optima (or) inflection-points.

Considering Newton-Raphson method, the update rule for root-finding is:

\begin{align}
 x_{i+1} = x_i - \dfrac{f(x_i)}{f'(x_i)}
\end{align}
 
Using the same Newton-Raphson method, the update rule for optimization is:
 
\begin{align}
x_{i+1} = x_i - \dfrac{f'(x_i)}{f''(x_i)}
\end{align}
 
This approach to optimization is called the \textbf{Newton's Method} in one dimension. It is worth highlighting that this is an open method for optimization. 
  
\begin{example} Determine a local optima of the function $f(x) = exp(x)-2x$ using the Newton-Raphson method
\begin{fullwidth}
  	\begin{verbatim}
%% Newton Raphson
	clear all;
	close all;
	x_start=100;
	f = @(x) exp(x)-2*x;
	f_dash = @(x) exp(x)-2;
	f_ddash = @(x) exp(x);
	tol=1e-8;
	tic;
	x_old=x_start;
	while (1) 
	x_new = x_old-f_dash(x_old)/f_ddash(x_old); %From similar triangles 
	%Convergence criterion
	if(abs((x_new-x_old)/x_new)<tol)
	break;
	end
	x_old=x_new;
	end
	disp('Minima found:') 
	disp(x_new)
	toc;
  	\end{verbatim}
\end{fullwidth}

\begin{margintable}
	\textbf{Output:}
	\vspace{-8pt}
	\begin{verbatim}
	Minima found: 0.6931
	Elapsed time is 0.006971 seconds.
	\end{verbatim}
\end{margintable}
\end{example}
	
In a similar manner, other root finding techniques can also be used for optimization to converge to local optima where the derivate is zero.
  
%\section{Matlab Functions for Optimization}
%The function \texttt{\textbf{fminbnd}} is designed to find the optimum of a single variable function within a bracket. The fminbnd function works as follows. It combines the slow, dependable golden-section search with the faster, but possibly unreliable, parabolic interpolation. It first attempts parabolic interpolation and keeps applying it as long as acceptable results are obtained. If not, it uses the golden-section search to get matters in hand.

%\begin{margintable}
%	\textbf{Syntax of \texttt{fminbnd}:}
%	\begin{verbatim}
%	[xmin, fval] = fminbnd(function,x1,x2)
%	\end{verbatim}
%\end{margintable}
%A simple representation of its syntax is \texttt{\textbf{fminbnd}(function,x1,x2)}, where \textit{x} and \textit{fval} are the location and value of the minimum, \textit{function} is the name of the function being evaluated, and ${x1}$ and ${x2}$ are the bounds of the interval being searched.

%\begin{margintable}
%	\textbf{Displaying iteration details:}
%	\begin{verbatim}
%	options = optimset('display','iter');
%	fminbnd(function,x1,x2,options)
%	\end{verbatim}
%\end{margintable}

%To display additional details about each iteration, optional parameters can be specified using \textit{optimset}. For example, we can display the calculation details as given to the right.\\
%\textbf{Example: } Find the minima of the unimodal function $f(x) = e^x-2x$ using the Matlab-function \textbf{fminbnd} with initial bounds of $x_1 = 0$, and $x_2 = 2$.\\
%\textbf{Solution: } 

%\begin{fullwidth}
%	\begin{Verbatim}
%>> f = @(x) exp(x)-2*x;	
%>> options = optimset('display','iter');
%>> fminbnd(f,0,2,options)
%	\end{Verbatim}
%\end{fullwidth}

%The above commands, when entered in the command-window, displays the following result:
%\begin{fullwidth}
%	\begin{Verbatim}
%Func-count     x          f(x)         Procedure
% 1         0.763932     0.618836        initial
% 2          1.23607     0.969917        golden
% 3         0.472136     0.659143        golden
% 4         0.677873     0.613938        parabolic
% 5         0.690883     0.613711        parabolic
% 6         0.692947     0.613706        parabolic
% 7         0.693154     0.613706        parabolic
% 8          0.69312     0.613706        parabolic
% 9         0.693187     0.613706        parabolic
% Optimization terminated:  the current x satisfies the termination 
% criteria using OPTIONS.TolX of 1e-04   ans = 0.6932
%	\end{Verbatim}
%\end{fullwidth}

%Thus, after three iterations, the method switches from golden to parabolic, and after nine iterations, the minimum is determined to a tolerance of 0.0001.
 
%For multi-dimensional optimization in Matlab, the function \texttt{\textbf{fminsearch}} can be used.\vspace{-12pt}
\subsection{Exercises}
\begin{enumerate}
    \item TODO
\end{enumerate}
\end{document}
